\documentclass[a4paper,10pt]{article}
\usepackage[francais]{babel}
\usepackage{url}

\selectlanguage{francais}
\input{letter-icube-uds.sty}

\selectlanguage{francais}
\begin{document}

\name{St\'ephane \textsc{Genaud}}
\fonction{Professeur des Universités}
\telephoneillkirch{+33 (0)3688 54542}
\email{genaud@unistra.fr}


\begin{letter}[Invitation au séminaire équipe ICPS, laboratoire Icube]
  {à}{Nestor Demeure, doctorant ENS, Paris-Saclay - CEA}

  Cher Nestor,

  je veux vous inviter  par la présente lettre à venir  présenter vos travaux de
  recherche lors  de notre équipe de  recherche ICPS du laboratoire  Icube, dans
  ses locaux situés 400 Blvd S. Brant, 67400 Illkirch, le 28 novembre 2018.

  La proposition que vous nous avez faite:\\

  \emph{Shaman : Une méthode directe et déterministe pour estimer l’erreur due à l’arithmétique flottante.}

  \begin{quote}
    Le résultat des calculs avec des nombres à virgule flottante est-il précis ?
    Est-ce qu'il  reste précis  lorsque la  taille du  problème augmente  ? Dans
    cette présentation nous introduisons  l'utilisation d'arithmétique par paire
    pour estimer la perte de précision  lié à l'arithmétique à virgule flottante
    et  fournissons une  implémentation  de référence,  la bibliothèque  Shaman.
    Notre approche permet l'estimation  de l'erreur introduite par l'utilisation
    de  l'arithmétique à  virgule  flottante en  tout point  du  code, elle  est
    déterministe  et, comme  l'illustrent plusieurs  exemples, compétitive  avec
    l'état de  l'art tout en étant  adaptée au calcul haute  performance grâce à
    son faible impact sur les temps de calcul.
\end{quote}

  a retenu notre attention.

  Ne disposant pas de budget propre pour ces déplacements, nous ne serons pas en
  mesure de vous indemniser du déplacement.

  Cordialement,

\end{letter}
\begin{flushright}
\includegraphics[width=.20\textwidth]{signgenaud.jpg}
\end{flushright}
\end{document}
